% Options for packages loaded elsewhere
% Options for packages loaded elsewhere
\PassOptionsToPackage{unicode}{hyperref}
\PassOptionsToPackage{hyphens}{url}
%
\documentclass[
  english,
  russian,
  12pt,
  a4paper,
  DIV=11,
  numbers=noendperiod]{scrreprt}
\usepackage{xcolor}
\usepackage{amsmath,amssymb}
\setcounter{secnumdepth}{5}
\usepackage{iftex}
\ifPDFTeX
  \usepackage[T1]{fontenc}
  \usepackage[utf8]{inputenc}
  \usepackage{textcomp} % provide euro and other symbols
\else % if luatex or xetex
  \usepackage{unicode-math} % this also loads fontspec
  \defaultfontfeatures{Scale=MatchLowercase}
  \defaultfontfeatures[\rmfamily]{Ligatures=TeX,Scale=1}
\fi
\usepackage{lmodern}
\ifPDFTeX\else
  % xetex/luatex font selection
\fi
% Use upquote if available, for straight quotes in verbatim environments
\IfFileExists{upquote.sty}{\usepackage{upquote}}{}
\IfFileExists{microtype.sty}{% use microtype if available
  \usepackage[]{microtype}
  \UseMicrotypeSet[protrusion]{basicmath} % disable protrusion for tt fonts
}{}
\usepackage{setspace}
% Make \paragraph and \subparagraph free-standing
\makeatletter
\ifx\paragraph\undefined\else
  \let\oldparagraph\paragraph
  \renewcommand{\paragraph}{
    \@ifstar
      \xxxParagraphStar
      \xxxParagraphNoStar
  }
  \newcommand{\xxxParagraphStar}[1]{\oldparagraph*{#1}\mbox{}}
  \newcommand{\xxxParagraphNoStar}[1]{\oldparagraph{#1}\mbox{}}
\fi
\ifx\subparagraph\undefined\else
  \let\oldsubparagraph\subparagraph
  \renewcommand{\subparagraph}{
    \@ifstar
      \xxxSubParagraphStar
      \xxxSubParagraphNoStar
  }
  \newcommand{\xxxSubParagraphStar}[1]{\oldsubparagraph*{#1}\mbox{}}
  \newcommand{\xxxSubParagraphNoStar}[1]{\oldsubparagraph{#1}\mbox{}}
\fi
\makeatother


\usepackage{longtable,booktabs,array}
\usepackage{calc} % for calculating minipage widths
% Correct order of tables after \paragraph or \subparagraph
\usepackage{etoolbox}
\makeatletter
\patchcmd\longtable{\par}{\if@noskipsec\mbox{}\fi\par}{}{}
\makeatother
% Allow footnotes in longtable head/foot
\IfFileExists{footnotehyper.sty}{\usepackage{footnotehyper}}{\usepackage{footnote}}
\makesavenoteenv{longtable}
\usepackage{graphicx}
\makeatletter
\newsavebox\pandoc@box
\newcommand*\pandocbounded[1]{% scales image to fit in text height/width
  \sbox\pandoc@box{#1}%
  \Gscale@div\@tempa{\textheight}{\dimexpr\ht\pandoc@box+\dp\pandoc@box\relax}%
  \Gscale@div\@tempb{\linewidth}{\wd\pandoc@box}%
  \ifdim\@tempb\p@<\@tempa\p@\let\@tempa\@tempb\fi% select the smaller of both
  \ifdim\@tempa\p@<\p@\scalebox{\@tempa}{\usebox\pandoc@box}%
  \else\usebox{\pandoc@box}%
  \fi%
}
% Set default figure placement to htbp
\def\fps@figure{htbp}
\makeatother



\ifLuaTeX
\usepackage[bidi=basic,provide=*]{babel}
\else
\usepackage[bidi=default,provide=*]{babel}
\fi
% get rid of language-specific shorthands (see #6817):
\let\LanguageShortHands\languageshorthands
\def\languageshorthands#1{}


\setlength{\emergencystretch}{3em} % prevent overfull lines

\providecommand{\tightlist}{%
  \setlength{\itemsep}{0pt}\setlength{\parskip}{0pt}}



 
\usepackage[style=gost-numeric,backend=biber,langhook=extras,autolang=other*]{biblatex}
\addbibresource{bib/cite.bib}

\usepackage[]{csquotes}

\usepackage{indentfirst}
\usepackage{float}
\floatplacement{figure}{H}
\IfFileExists{plex-otf.sty}{
  %% Full TeXlive
  \usepackage[
    % math,
    RM={Scale=0.94},SS={Scale=0.94},SScon={Scale=0.94},TT={Scale=MatchLowercase,FakeStretch=0.9},
    DefaultFeatures={Ligatures=Common}
  ]{plex-otf}
  % \usepackage[Scale=MatchUppercase]{juliamono}
}{
  %% TinyTeX
  \usepackage{libertine}
}
\KOMAoption{captions}{tableheading}
\makeatletter
\@ifpackageloaded{caption}{}{\usepackage{caption}}
\AtBeginDocument{%
\ifdefined\contentsname
  \renewcommand*\contentsname{Содержание}
\else
  \newcommand\contentsname{Содержание}
\fi
\ifdefined\listfigurename
  \renewcommand*\listfigurename{Список иллюстраций}
\else
  \newcommand\listfigurename{Список иллюстраций}
\fi
\ifdefined\listtablename
  \renewcommand*\listtablename{Список таблиц}
\else
  \newcommand\listtablename{Список таблиц}
\fi
\ifdefined\figurename
  \renewcommand*\figurename{Рисунок}
\else
  \newcommand\figurename{Рисунок}
\fi
\ifdefined\tablename
  \renewcommand*\tablename{Таблица}
\else
  \newcommand\tablename{Таблица}
\fi
}
\@ifpackageloaded{float}{}{\usepackage{float}}
\floatstyle{ruled}
\@ifundefined{c@chapter}{\newfloat{codelisting}{h}{lop}}{\newfloat{codelisting}{h}{lop}[chapter]}
\floatname{codelisting}{Список}
\newcommand*\listoflistings{\listof{codelisting}{Листинги}}
\makeatother
\makeatletter
\makeatother
\makeatletter
\@ifpackageloaded{caption}{}{\usepackage{caption}}
\@ifpackageloaded{subcaption}{}{\usepackage{subcaption}}
\makeatother
\usepackage{bookmark}
\IfFileExists{xurl.sty}{\usepackage{xurl}}{} % add URL line breaks if available
\urlstyle{same}
\hypersetup{
  pdftitle={отчёт по лабораторной работе №},
  pdfauthor={Кхари Жекка Кализая арсе},
  pdflang={ru-RU},
  hidelinks,
  pdfcreator={LaTeX via pandoc}}


\title{отчёт по лабораторной работе №}
\usepackage{etoolbox}
\makeatletter
\providecommand{\subtitle}[1]{% add subtitle to \maketitle
  \apptocmd{\@title}{\par {\large #1 \par}}{}{}
}
\makeatother
\subtitle{Первоначальное конфигурирование сети}
\author{Кхари Жекка Кализая арсе}
\date{}
\begin{document}
\maketitle

\renewcommand*\contentsname{Содержание}
{
\setcounter{tocdepth}{1}
\tableofcontents
}
\listoffigures
\listoftables

\setstretch{1.5}
\chapter{Цель
работы}\label{ux446ux435ux43bux44c-ux440ux430ux431ux43eux442ux44b}

Провести подготовительную работу по первоначальной настройке
коммутаторов сети.

\chapter{Задание}\label{ux437ux430ux434ux430ux43dux438ux435}

Требуется сделать первоначальную настройку коммутаторов сети,
представленной на схеме L1 (см. рис. 3.1 из раздела 3.3). Под
первоначальной настройкой понимается указание имени устройства, его
IP-адреса, настройка доступа по паролю к виртуальным терминалам и
консоли, настройка удалённого доступа к устройству по ssh. При
выполнении работы необходимо учитывать соглашение об именовании (см.
раздел 2.5).

\chapter{Выполнение лабораторной
работы}\label{ux432ux44bux43fux43eux43bux43dux435ux43dux438ux435-ux43bux430ux431ux43eux440ux430ux442ux43eux440ux43dux43eux439-ux440ux430ux431ux43eux442ux44b}

\section{создание
топологии}\label{ux441ux43eux437ux434ux430ux43dux438ux435-ux442ux43eux43fux43eux43bux43eux433ux438ux438}

Сначала былы расположенны все коммутаторы в рабочей области
(рис.~\ref{fig-001}), здесь былы использованые один коммутатор 2960 и
четыре коммутатора 2950.

\begin{figure}

\centering{

\includegraphics[width=0.7\linewidth,height=\textheight,keepaspectratio]{image/fig00001.png}

}

\caption{\label{fig-001}коммутаторы}

\end{figure}%

Потом коммутаторы были соединенные проводами следуя таблицу быдно в
предыдущей лабораторной работе (рис.~\ref{fig-002}).

\begin{figure}

\centering{

\includegraphics[width=0.7\linewidth,height=\textheight,keepaspectratio]{image/fig00002.png}

}

\caption{\label{fig-002}соединение коммутаторов}

\end{figure}%

Дальше были расположенны обычные компьютеры и серверы
(рис.~\ref{fig-003}).

\begin{figure}

\centering{

\includegraphics[width=0.7\linewidth,height=\textheight,keepaspectratio]{image/fig00003.png}

}

\caption{\label{fig-003}компьютеры и серверы}

\end{figure}%

Затем были соединенны с помощью проводов (рис.~\ref{fig-004}).

\begin{figure}

\centering{

\includegraphics[width=0.7\linewidth,height=\textheight,keepaspectratio]{image/fig00004.png}

}

\caption{\label{fig-004}соединение компьютеров и серверов}

\end{figure}%

\section{настройка
коммутаторы}\label{ux43dux430ux441ux442ux440ux43eux439ux43aux430-ux43aux43eux43cux43cux443ux442ux430ux442ux43eux440ux44b}

Сначала я щелкал коммутатор 1 чтобы открывать окно и там я щелкал CLI.
там я смог смотреть терминал, с которым я мог настроить коммутатор.

Сначала я изменил hostname на msk-donskaya-qscalizaya-sw-1, потом я
начал настройку интерфейса VLAN2. я включил его и дал IP-адрес
10.128.1.2 и маску подсети 255.255.255.0 (этот IP-адрес будет измениться
в каждом коммутаторе следуя таблицу IP-адресов) дальше я указал по
умолчанию gateway 10.128.1.1. Дальше я настроил удаленный доступ

Затем настраиваются линии VTY (Virtual Teletype), чтобы обеспечить
возможность одновременного подключения от 1 до 5 сессий (0--4). Далее
задаётся пароль, который будет использоваться для этих сессий (cisco), и
с помощью команды login принудительно активируется требование ввода
пароля.

После этого настраивается физический доступ к коммутатору через консоль.
снова устанавливается пароль (cisco) и принудительно включается его
использование.

В завершение задаётся пароль (cisco) для привилегированного режима и
активируется служба шифрования паролей. (рис.~\ref{fig-005}) .

\begin{figure}

\centering{

\includegraphics[width=0.7\linewidth,height=\textheight,keepaspectratio]{image/fig00005.png}

}

\caption{\label{fig-005}настройка VLAN2 и паролей}

\end{figure}%

Затем был настроен пользователь для подключения по SSH с именем
пользователя admin и паролем cisco, которому был назначен базовый
уровень привилегий 1. После этого было задано доменное имя
donskaya.rudn.edu. Также была создана RSA-ключевая пара с размером
модуля 800 бит.

Далее я вернулся к настройке линий виртуального терминала (VTY), чтобы
принудительно разрешить только подключение по SSH. После завершения
конфигурации я вышел из режима настройки и сохранил изменения в памяти
устройства. (рис.~\ref{fig-006}).

\begin{figure}

\centering{

\includegraphics[width=0.7\linewidth,height=\textheight,keepaspectratio]{image/fig00006.png}

}

\caption{\label{fig-006}настройка пользователя и протокола SSH для
vitual teletype}

\end{figure}%

Потом я сделал тоже самое для всех остальных коммутаторов. (с
рис.~\ref{fig-007}).

ip-адрес VLAN2:

\begin{itemize}
\tightlist
\item
  msk-donskaya-qscalizaya-sw-2: 10.128.1.3 (рис.~\ref{fig-007}).
\item
  msk-donskaya-qscalizaya-sw-3: 10.128.1.4 (рис.~\ref{fig-008}).
\item
  msk-donskaya-qscalizaya-sw-4: 10.128.1.5 (рис.~\ref{fig-009}).
\item
  msk-pavlovskaya-qscalizaya-sw-1: 10.128.1.6 (рис.~\ref{fig-010}).
\end{itemize}

\begin{figure}

\centering{

\includegraphics[width=0.7\linewidth,height=\textheight,keepaspectratio]{image/fig00007.png}

}

\caption{\label{fig-007}настройка коммутатора
msk-donskaya-qscalizaya-sw-2}

\end{figure}%

\begin{figure}

\centering{

\includegraphics[width=0.7\linewidth,height=\textheight,keepaspectratio]{image/fig00008.png}

}

\caption{\label{fig-008}настройка коммутатора
msk-donskaya-qscalizaya-sw-3}

\end{figure}%

\begin{figure}

\centering{

\includegraphics[width=0.7\linewidth,height=\textheight,keepaspectratio]{image/fig00009.png}

}

\caption{\label{fig-009}настройка коммутатора
msk-donskaya-qscalizaya-sw-4}

\end{figure}%

\begin{figure}

\centering{

\includegraphics[width=0.7\linewidth,height=\textheight,keepaspectratio]{image/fig00010.png}

}

\caption{\label{fig-010}настройка коммутатора
msk-pavlovskaya-qscalizaya-sw-1}

\end{figure}%

Дальше я проверил работоспособность системы

\begin{figure}

\centering{

\includegraphics[width=0.7\linewidth,height=\textheight,keepaspectratio]{image/fig00011.png}

}

\caption{\label{fig-011}проверка системы}

\end{figure}%

\begin{figure}

\centering{

\includegraphics[width=0.7\linewidth,height=\textheight,keepaspectratio]{image/fig00012.png}

}

\caption{\label{fig-012}проверка системы}

\end{figure}%

\chapter{Выводы}\label{ux432ux44bux432ux43eux434ux44b}

В этой лабораторной работе я смог смотреть как можно настроить
коммутаторы чтобы указывать IP-адрес VLAN2 и также как настроить Virtual
teletype пароли и протокол SSH

\chapter*{Список
литературы}\label{ux441ux43fux438ux441ux43eux43a-ux43bux438ux442ux435ux440ux430ux442ux443ux440ux44b}
\addcontentsline{toc}{chapter}{Список литературы}

\printbibliography[heading=none]





\end{document}
